Elemento opcional (Guia UECE 2026, 2.2.2.1.3), usado **somente** quando um erro é percebido depois que o trabalho já foi impresso ou depositado. Não é para ser preenchido durante a escrita normal do TCC -- se você não tiver nenhuma correção a registrar, apague o conteúdo abaixo e remova a linha \texttt{\textbackslash imprimirerrata\{...\}} do documento.tex.

Quando for necessária, a errata é composta por dois elementos, nesta ordem: (1) a referência completa do próprio trabalho, no padrão ABNT NBR 6023, e (2) uma tabela indicando onde está o erro e qual é a correção, com as colunas Folha, Linha, Onde se lê e Leia-se.

Substitua os dados abaixo pelos do seu próprio trabalho e pelas correções identificadas:

\noindent SOBRENOME, Nome do autor. \textbf{Título do trabalho}: subtítulo (se houver). Ano. Número de folhas f. Trabalho de Conclusão de Curso (Graduação em [curso]) -- Centro de Ciências e Tecnologia, Universidade Estadual do Ceará, Fortaleza, ano.

\begin{table}[ht!]
	\center
	\begin{tabular}{|p{1.4cm}|p{1cm}|p{3cm}|p{3cm}|}
		\hline
		\textbf{Folha} & \textbf{Linha} & \textbf{Onde se lê} &
		\textbf{Leia-se}\\
		\hline
		00 & 00 & palavra ou trecho incorreto & palavra ou trecho correto\\
		\hline
	\end{tabular}
\end{table}
